\documentclass[12pt,a4paper]{article}
\usepackage[utf8]{inputenc}
\usepackage[T1]{fontenc}
\usepackage[spanish]{babel}
\usepackage{lmodern}
\usepackage{geometry}
\usepackage{hyperref}
\usepackage{enumitem}
\usepackage{parskip}
\usepackage{longtable}
\usepackage{array}
\usepackage{xcolor}

\geometry{margin=2.2cm}
\hypersetup{
  colorlinks=true,
  linkcolor=black,
  urlcolor=blue,
  pdftitle={Plan de Trabajo Profesional - The Daily Scribble},
  pdfauthor={Harry \& EPN}
}

\setlist[itemize]{leftmargin=*}
\setlist[enumerate]{leftmargin=*}

\title{Plan de Trabajo Profesional\\\large The Daily Scribble}
\author{Proyecto Personal con Enfoque Profesional}
\date{30 de marzo de 2026}

\begin{document}

\maketitle
\tableofcontents
\newpage

\section{Resumen Ejecutivo}
Este documento define un plan de trabajo integral para evolucionar \textbf{The Daily Scribble} como un proyecto de software serio, mantenible y escalable. El objetivo es profesionalizar el producto sin perder su esencia: dibujo diario rapido, creatividad sin friccion y galeria comunitaria efimera.

La estrategia se apoya en cinco pilares:
\begin{itemize}
  \item Producto: alcance claro, prioridades y roadmap por etapas.
  \item Ingenieria: arquitectura modular, estandares de codigo y calidad.
  \item Operacion: procesos de versionado, CI/CD y manejo de incidencias.
  \item Datos y metricas: monitoreo de uso, calidad y salud del sistema.
  \item Documentacion: especificaciones vivas para reducir deuda tecnica.
\end{itemize}

\section{Estado Actual (AS-IS) y Supuestos Confirmados}
A continuacion se consolida el estado funcional actual y las decisiones de producto expresadas por el equipo:

\subsection{Comportamiento actual y restricciones}
\begin{itemize}
  \item El canvas debe ser responsivo al tamano de pantalla en PC (no movil por ahora).
  \item Si el usuario cambia el tamano de ventana, el canvas se reescala.
  \item No existe soporte de presion de lapiz.
  \item Limite de historial: 4 estados (objetivo de producto para siguiente fase).
  \item Gotero solo toma color del lienzo interno.
  \item No existen capas (layers).
  \item Palabra diaria proveniente de una lista de 100 palabras.
  \item La palabra del dia es la misma para todos los usuarios.
  \item Deben existir alertas visuales y contador visual.
  \item El temporizador no se puede pausar.
  \item Si el usuario refresca (F5), la sesion debe continuar.
  \item Un usuario puede hacer mas de un dibujo si elige distinta duracion.
  \item Galeria muestra un maximo de 30 dibujos por ahora.
  \item No se pueden borrar dibujos.
  \item Al clickear un dibujo en galeria, se expande y admite like/dislike.
  \item Dibujos con mas likes tienen prioridad de visualizacion.
  \item La galeria se limpia cada dia.
  \item No se prioriza usabilidad movil en esta etapa.
  \item No hay modo oscuro en esta etapa, pero se considera para roadmap.
  \item No se requiere sistema de duelos ni logros.
\end{itemize}

\subsection{Riesgos inmediatos del estado actual}
\begin{itemize}
  \item Deuda tecnica por crecimiento de funcionalidades en un unico archivo JS.
  \item Escalabilidad limitada de galeria al crecer volumen de dibujos.
  \item Ausencia de backend real para ranking global robusto.
  \item Riesgo de inconsistencias al recuperar estado tras refresco.
\end{itemize}

\section{Vision, Objetivos y Principios}
\subsection{Vision}
Construir una plataforma creativa diaria, simple y confiable, con experiencia cuidada en escritorio y bases tecnicas para escalar a comunidad real.

\subsection{Objetivos estrategicos (12 meses)}
\begin{enumerate}
  \item Entregar experiencia de dibujo de escritorio estable, fluida y predecible.
  \item Consolidar galeria diaria con orden por relevancia social (likes/dislikes).
  \item Preparar arquitectura para trafico creciente sin reescritura total.
  \item Institucionalizar practicas profesionales: QA, CI/CD, trazabilidad y documentacion.
\end{enumerate}

\subsection{Principios de producto e ingenieria}
\begin{itemize}
  \item Simplicidad primero, sin sacrificar robustez.
  \item Decisiones guiadas por metricas y evidencia.
  \item Seguridad y privacidad por defecto.
  \item Documentacion como parte del entregable, no como anexo.
  \item Evolucion incremental: pequenos pasos con valor real.
\end{itemize}

\section{Alcance}
\subsection{Alcance en esta hoja de ruta}
\begin{itemize}
  \item Desktop-first responsivo.
  \item Persistencia de sesion frente a refresco F5.
  \item Galeria diaria con ranking por votos.
  \item Limite de 30 items visibles y estrategia de crecimiento progresiva.
  \item Gobierno de codigo y flujo profesional de entregas.
\end{itemize}

\subsection{Fuera de alcance inicial}
\begin{itemize}
  \item Aplicacion movil nativa o web optimizada para moviles.
  \item Modo oscuro productivo completo.
  \item Mecanicas de juego (duelos, logros, torneos).
\end{itemize}

\section{Arquitectura Objetivo (TO-BE)}
\subsection{Modelo tecnico recomendado}
\begin{itemize}
  \item Frontend modular por dominios:
  \begin{itemize}
    \item modulo UI y vistas,
    \item modulo motor de dibujo,
    \item modulo temporizador,
    \item modulo persistencia,
    \item modulo galeria,
    \item modulo votos/ranking.
  \end{itemize}
  \item Estado de aplicacion centralizado (store liviano con eventos).
  \item Contratos de datos tipados y versionados (aunque sea JS, con esquema JSON).
  \item Adaptador de almacenamiento para permitir migracion de localStorage a backend.
\end{itemize}

\subsection{Persistencia de sesion}
Para cumplir continuidad tras F5, se define una entidad \texttt{drawingSession} con:
\begin{itemize}
  \item identificador de sesion,
  \item duracion elegida,
  \item timestamp de inicio,
  \item segundos restantes,
  \item snapshot del canvas,
  \item palabra diaria,
  \item estado de bloqueo por fin de tiempo.
\end{itemize}

Regla de recuperacion:
\begin{enumerate}
  \item Al iniciar app, leer sesion activa.
  \item Recalcular tiempo restante con reloj real (no confiar solo en contador en memoria).
  \item Si tiempo restante mayor a 0, restaurar lienzo y continuar.
  \item Si tiempo restante menor o igual a 0, bloquear edicion y mostrar alerta visual.
\end{enumerate}

\section{Roadmap por Etapas (Plan Maestro)}
\subsection{Etapa 0: Gobierno del proyecto y base profesional}
\textbf{Duracion sugerida: 1 semana}

\textbf{Objetivo}: establecer cimientos de gestion y calidad.

\textbf{Entregables}:
\begin{itemize}
  \item Documento de vision y alcance versionado.
  \item Convenciones de codigo y estilo.
  \item Politica de ramas y estrategia de commits.
  \item Plantillas de issue, bug report y pull request.
  \item Definition of Ready (DoR) y Definition of Done (DoD).
\end{itemize}

\textbf{Criterio de salida}: todo cambio nuevo entra por PR con checklist de calidad.

\subsection{Etapa 1: Refactor tecnico y modularizacion}
\textbf{Duracion sugerida: 2 semanas}

\textbf{Objetivo}: reducir deuda tecnica y preparar crecimiento.

\textbf{Entregables}:
\begin{itemize}
  \item Separacion de \texttt{app.js} en modulos.
  \item Capa de estado centralizada.
  \item Tests base unitarios para utilidades criticas.
  \item Logging funcional para eventos clave.
\end{itemize}

\textbf{Criterio de salida}: reduccion de complejidad ciclomatica y cobertura minima acordada.

\subsection{Etapa 2: Canvas desktop responsivo serio}
\textbf{Duracion sugerida: 2 semanas}

\textbf{Objetivo}: experiencia de dibujo consistente en escritorio al redimensionar ventana.

\textbf{Entregables}:
\begin{itemize}
  \item Algoritmo de resize preservando contenido del canvas.
  \item Escalado proporcional segun area visible.
  \item Limite de historial configurado en 4 estados.
  \item Validacion de gotero restringido al lienzo.
  \item Suite de pruebas de regresion visual minima.
\end{itemize}

\textbf{Criterio de salida}: sin perdida de dibujo en resize en escenarios definidos de QA.

\subsection{Etapa 3: Temporizador robusto y continuidad F5}
\textbf{Duracion sugerida: 2 semanas}

\textbf{Objetivo}: hacer el temporizador confiable y persistente.

\textbf{Entregables}:
\begin{itemize}
  \item Contador visual con estados (normal, advertencia, critico).
  \item Alertas visuales por umbrales (por ejemplo 60s y 10s).
  \item Sin opcion de pausa.
  \item Continuidad de sesion tras refresco F5.
  \item Bloqueo de dibujo al finalizar tiempo con feedback claro.
\end{itemize}

\textbf{Criterio de salida}: tasa de fallos de recuperacion de sesion menor al 1\% en pruebas.

\subsection{Etapa 4: Galeria diaria v1 con ranking social}
\textbf{Duracion sugerida: 3 semanas}

\textbf{Objetivo}: consolidar exploracion de obras y priorizacion por votos.

\textbf{Entregables}:
\begin{itemize}
  \item Listado con tope de 30 obras visibles.
  \item Modal o vista expandida al clickear una obra.
  \item Acciones like/dislike por obra.
  \item Regla de orden: score de votos y desempate por recencia.
  \item Restriccion de no borrado de dibujos.
\end{itemize}

\textbf{Criterio de salida}: ranking estable y UX clara en interaccion de voto.

\subsection{Etapa 5: Datos diarios y palabra global}
\textbf{Duracion sugerida: 1 semana}

\textbf{Objetivo}: formalizar reglas diarias de contenido.

\textbf{Entregables}:
\begin{itemize}
  \item Lista curada de 100 palabras versionada.
  \item Regla deterministica para palabra del dia global.
  \item Limpieza automatica de galeria por cambio de fecha.
  \item Pruebas de borde en cambio de dia y zonas horarias.
\end{itemize}

\textbf{Criterio de salida}: todos los usuarios observan la misma palabra del dia bajo la misma fecha de referencia.

\subsection{Etapa 6: Observabilidad y operacion}
\textbf{Duracion sugerida: 2 semanas}

\textbf{Objetivo}: operar con estandares de producto real.

\textbf{Entregables}:
\begin{itemize}
  \item Telemetria de eventos funcionales (sin datos sensibles).
  \item Dashboard de metricas clave.
  \item Registro de errores con trazabilidad.
  \item Flujo de release con versionado semantico.
\end{itemize}

\textbf{Criterio de salida}: capacidad de detectar, diagnosticar y priorizar fallos con evidencia.

\subsection{Etapa 7: Preparacion para escalabilidad}
\textbf{Duracion sugerida: 3 semanas}

\textbf{Objetivo}: disenar transicion hacia backend para volumen alto.

\textbf{Entregables}:
\begin{itemize}
  \item Documento de arquitectura de referencia (API, almacenamiento, cache).
  \item Estrategia de paginacion de galeria y top-N por score.
  \item Modelo de votos resistente a abuso basico.
  \item Plan de migracion progresiva desde localStorage.
\end{itemize}

\textbf{Criterio de salida}: blueprint tecnico listo para implementacion sin incertidumbre mayor.

\subsection{Etapa 8: Backlog evolutivo (no prioritario inicial)}
\textbf{Duracion sugerida: continua}

\textbf{Objetivo}: incorporar mejoras diferidas (movil, modo oscuro, etc.) de forma controlada.

\textbf{Entregables}: RFCs de nuevas capacidades y prototipos validados.

\section{Backlog Priorizado Inicial}
\subsection{Epicas}
\begin{enumerate}
  \item E1: Motor de dibujo robusto desktop.
  \item E2: Temporizador persistente y confiable.
  \item E3: Galeria social diaria con ranking.
  \item E4: Calidad, procesos y observabilidad.
\end{enumerate}

\subsection{Historias de usuario (muestra)}
\begin{itemize}
  \item HU-01: Como usuario desktop, quiero que el canvas se adapte al tamano de mi ventana sin perder mi dibujo.
  \item HU-02: Como usuario, quiero que al refrescar F5 mi sesion siga donde quedo.
  \item HU-03: Como usuario, quiero votar un dibujo con like/dislike al expandirlo.
  \item HU-04: Como usuario, quiero ver primero los dibujos mejor valorados.
  \item HU-05: Como usuario, quiero hacer mas de un dibujo si cambio la duracion.
  \item HU-06: Como usuario, quiero alertas visuales cuando el tiempo este por terminar.
\end{itemize}

\subsection{Criterios de aceptacion (ejemplos)}
\begin{itemize}
  \item CA-01: Al redimensionar ventana, el trazo previo permanece visible y proporcional.
  \item CA-02: Al refrescar, el contador se reconstruye con diferencia de tiempo real.
  \item CA-03: El sistema no muestra boton de pausa ni permite pausar por atajos.
  \item CA-04: Galeria no muestra mas de 30 elementos en la vista principal.
  \item CA-05: Un click en miniatura abre vista expandida y controles de voto.
\end{itemize}

\section{Plan de Calidad}
\subsection{Estrategia de pruebas}
\begin{itemize}
  \item Unitarias: utilidades de tiempo, ranking y seleccion de palabra.
  \item Integracion: flujo sesion completa (inicio, dibujo, fin, guardado).
  \item End-to-end: casos criticos en navegador desktop.
  \item Regresion visual: canvas y componentes de galeria.
  \item Pruebas manuales guiadas para UX de dibujo.
\end{itemize}

\subsection{Matriz minima de pruebas}
\begin{itemize}
  \item Navegadores objetivo: Chrome, Edge, Firefox (desktop).
  \item Resoluciones: 1366x768, 1600x900, 1920x1080, 2560x1440.
  \item Escenarios: resize agresivo, F5 durante dibujo, fin de tiempo, votos en galeria.
\end{itemize}

\subsection{Definicion de Done por historia}
\begin{itemize}
  \item Criterios de aceptacion validados.
  \item Sin errores criticos abiertos.
  \item Cobertura de pruebas actualizada.
  \item Documentacion tecnica y funcional actualizada.
  \item Evidencia de QA adjunta en PR.
\end{itemize}

\section{Plan de Documentacion}
\subsection{Documentos obligatorios}
\begin{itemize}
  \item PRD (Product Requirements Document).
  \item ADR (Architecture Decision Records).
  \item Especificacion de API/eventos internos.
  \item Guia de contribucion y estilo.
  \item Runbook operativo (incidentes y despliegues).
  \item Plan de pruebas y reporte de calidad por release.
\end{itemize}

\subsection{Cadencia de actualizacion}
\begin{itemize}
  \item PRD: revision mensual o ante cambios de alcance.
  \item ADR: uno por decision de arquitectura relevante.
  \item Runbook: actualizacion por incidente significativo.
\end{itemize}

\section{Plan de Versionado, CI/CD y Entrega}
\subsection{Flujo de ramas recomendado}
\begin{itemize}
  \item \texttt{main}: estable y desplegable.
  \item \texttt{develop} (opcional si se requiere): integracion.
  \item \texttt{feature/*}: trabajo por historia.
  \item \texttt{hotfix/*}: correcciones urgentes.
\end{itemize}

\subsection{Politica de commits}
\begin{itemize}
  \item Convencion semantica (ejemplo: feat, fix, refactor, docs, test, chore).
  \item Un cambio logico por commit.
  \item Mensajes trazables a issue/historia.
\end{itemize}

\subsection{Pipeline CI minimo}
\begin{enumerate}
  \item Lint y formato.
  \item Tests unitarios.
  \item Build/verificacion estatica.
  \item Smoke test de flujo principal.
  \item Publicacion automatizada a entorno de preview.
\end{enumerate}

\section{Seguridad, Privacidad y Abuso}
\begin{itemize}
  \item Minimizar datos personales almacenados.
  \item Validar entradas de texto para firma.
  \item Limitar spam de votos por huella local en etapa inicial.
  \item Preparar controles anti abuso para etapa backend.
\end{itemize}

\section{Matriz de Riesgos}
\begin{longtable}{|p{0.23\textwidth}|p{0.14\textwidth}|p{0.14\textwidth}|p{0.39\textwidth}|}
\hline
\textbf{Riesgo} & \textbf{Prob.} & \textbf{Impacto} & \textbf{Mitigacion} \\\hline
Perdida de dibujo en resize & Media & Alto & Implementar resize con canvas auxiliar y pruebas de regresion visual. \\\hline
Sesion inconsistente tras F5 & Media & Alto & Persistencia con timestamps y recomputo deterministico del tiempo. \\\hline
Ranking manipulable & Alta & Medio & Estrategia anti abuso inicial y migracion a backend con controles. \\\hline
Crecimiento de deuda tecnica & Alta & Alto & Modularizacion temprana, ADRs y code review estricto. \\\hline
Bajo rendimiento en galeria & Media & Medio & Limite de 30 visibles, lazy rendering y plan de paginacion. \\\hline
\end{longtable}

\section{KPI y Metricas de Exito}
\begin{itemize}
  \item Estabilidad: tasa de errores JS por sesion.
  \item Rendimiento: tiempo de carga inicial y tiempo de respuesta de interacciones.
  \item Retencion: usuarios que vuelven al dia siguiente.
  \item Conversion creativa: porcentaje de sesiones iniciadas que terminan en dibujo compartido.
  \item Calidad de experiencia: tasa de abandono durante sesion.
  \item Salud tecnica: cobertura de pruebas y deuda tecnica abierta.
\end{itemize}

\section{Plan de 30/60/90 dias}
\subsection{Primeros 30 dias}
\begin{itemize}
  \item Etapa 0 y Etapa 1 completadas.
  \item Base modular, convenciones y pipeline basico activos.
  \item Backlog refinado y priorizado para 2 sprints siguientes.
\end{itemize}

\subsection{Dias 31 a 60}
\begin{itemize}
  \item Etapa 2 y Etapa 3 completadas.
  \item Canvas responsivo desktop estable.
  \item Temporizador persistente y alertas visuales listas.
\end{itemize}

\subsection{Dias 61 a 90}
\begin{itemize}
  \item Etapa 4 y Etapa 5 completadas.
  \item Galeria con expander y votos en produccion inicial.
  \item Palabra diaria global de lista de 100 y limpieza diaria validada.
\end{itemize}

\section{Organizacion y Rituales de Trabajo}
\subsection{Roles minimos (aunque sea proyecto personal)}
\begin{itemize}
  \item Product Owner: prioriza alcance y acepta entregables.
  \item Tech Lead: define arquitectura y estandares.
  \item QA Owner: asegura cobertura y evidencia de calidad.
\end{itemize}

En proyecto personal, una misma persona puede cumplir todos los roles, pero se recomienda separar momentos de trabajo y checklist por rol para evitar sesgos.

\subsection{Rituales recomendados}
\begin{itemize}
  \item Planificacion semanal (1 hora).
  \item Revision de avance y riesgos (30 min, dos veces por semana).
  \item Demo interna por sprint.
  \item Retrospectiva corta para mejoras de proceso.
\end{itemize}

\section{Plan para capacidades futuras (movil y modo oscuro)}
Estas capacidades se consideran \textbf{futuras} y no bloquean el objetivo actual.

\subsection{Movil}
\begin{itemize}
  \item Definir breakpoints y simplificacion de panel de herramientas.
  \item Replantear interacciones de precision para tactil.
  \item Evaluar performance en dispositivos de gama media.
\end{itemize}

\subsection{Modo oscuro}
\begin{itemize}
  \item Diseñar sistema de tokens de color.
  \item Implementar selector de tema y preferencia persistente.
  \item Validar contraste y accesibilidad (WCAG).
\end{itemize}

\section{Checklist de Implementacion Inmediata}
\begin{enumerate}
  \item Crear tablero de trabajo (issues por epica e historias).
  \item Definir DoR/DoD y plantillas de PR.
  \item Ejecutar refactor modular de codigo actual.
  \item Implementar resize robusto del canvas y limite de historial en 4.
  \item Implementar persistencia de temporizador para continuidad F5.
  \item Construir modal de expansion en galeria con like/dislike.
  \item Ordenar galeria por score y limitar vista a 30.
  \item Configurar metricas y registro de errores.
\end{enumerate}

\section{Conclusion}
El proyecto tiene una base creativa fuerte y una direccion clara. Con este plan por etapas, \textbf{The Daily Scribble} puede evolucionar desde MVP personal hacia una plataforma tecnicamente madura, con entregas predecibles, calidad controlada y capacidad real de crecimiento.

La clave sera mantener disciplina de producto, trazabilidad tecnica y foco en experiencia de usuario en desktop, mientras se prepara la arquitectura para la siguiente fase de escalamiento.

\end{document}
